% GENERATED FILE. Edit canonical lessons, then run npm run generate:books.
% canonical-source-hash: fnv1a64:23ef06740174dcf5
% canonical-lessons: JA-C254-to-omoimasu, JA-C254-kara-node, JA-C254-iidasu, JA-W254-da-kanji, JA-C254-iken, JA-R254-what-i-think-and-why

\chapter{What I Think, and Why}
\label{ch:ja-what-i-think-and-why}
\hlchaptermodality{254}

\begin{chapteropening}
\textbf{By the end of this chapter:} \emph{I can say what I think with a plain clause and to omoimasu, ask dō omoimasu ka, give a reason with kara or node, and say who suggested something (iidashimashita).}

Complete JA-R254-what-i-think-and-why: say what you think and why, and ask someone else's opinion.
\end{chapteropening}

\section[\texorpdfstring{\ja{と} \ja{おもいます} (to omoimasu)}{と おもいます (to omoimasu)}]{\ja{と} \ja{おもいます} (to omoimasu) — I think that …}

\label{lesson:JA-C254-to-omoimasu}



\begin{quote}
Say \textquotedblleft{}to think\textquotedblright{}, and \textquotedblleft{}I think\textquotedblright{}, politely. (\emph{omou}, \emph{omoimasu}.) Then say, plainly, \textquotedblleft{}it will rain\textquotedblright{} and \textquotedblleft{}he won't come\textquotedblright{}. (\emph{ame ga furu}, \emph{konai}.) Those plain forms are about to get their second job.
\end{quote}

\subsection*{You'll want to know: a plain clause, then \ja{と}}
\begin{quote}
\textbf{\ja{いいと} \ja{おもいます。}}
\end{quote}

\begin{quote}
\emph{ii to omoimasu.} — \textquotedblleft{}I think it is good.\textquotedblright{}
\end{quote}

What you think goes first, in the plain form; \textbf{\ja{と}} (\emph{to}) closes it, and \textbf{\ja{おもいます}} (\emph{omoimasu}) keeps the whole sentence polite.

\begin{quote}
\textbf{\ja{あした} \ja{あめが} \ja{ふると} \ja{おもいます。}}
\end{quote}

\begin{quote}
\emph{ashita ame ga furu to omoimasu.} — \textquotedblleft{}I think it will rain tomorrow.\textquotedblright{}
\end{quote}

The plain \textquotedblleft{}not\textquotedblright{} and the plain past go before \textbf{\ja{と}} the same way:

\begin{quote}
\textbf{\ja{ともだちは} \ja{こないと} \ja{おもいます。}}
\end{quote}

\begin{quote}
\emph{tomodachi wa konai to omoimasu.} — \textquotedblleft{}I don't think my friend will come.\textquotedblright{}
\end{quote}

\begin{quote}
\textbf{\ja{ちちは} \ja{もう} \ja{かえったと} \ja{おもいます。}}
\end{quote}

\begin{quote}
\emph{chichi wa mō kaetta to omoimasu.} — \textquotedblleft{}I think my father has gone home already.\textquotedblright{}
\end{quote}

\subsection*{You'll want to know: \ja{どう} \ja{おもいますか}}
To ask what someone thinks, \textbf{\ja{どう}} (\emph{dō}), \textquotedblleft{}how\textquotedblright{}, goes where the opinion would be:

\begin{quote}
\textbf{\ja{どう} \ja{おもいますか。}}
\end{quote}

\begin{quote}
\emph{dō omoimasu ka.} — \textquotedblleft{}What do you think?\textquotedblright{}
\end{quote}

\begin{quote}
\textbf{\ja{この} \ja{えいがを} \ja{どう} \ja{おもいますか。}}
\end{quote}

\begin{quote}
\emph{kono eiga o dō omoimasu ka.} — \textquotedblleft{}What do you think of this film?\textquotedblright{}
\end{quote}

\begin{quote}
\textbf{\ja{おもしろいと} \ja{おもいます。}}
\end{quote}

\begin{quote}
\emph{omoshiroi to omoimasu.} — \textquotedblleft{}I think it is interesting.\textquotedblright{}
\end{quote}

\subsection*{Your turn}
Say these aloud:

\begin{itemize}
\raggedright
  \item \textquotedblleft{}I think it is interesting\textquotedblright{} — \emph{omoshiroi to omoimasu}
  \item \textquotedblleft{}I think it is good\textquotedblright{} — \emph{ii to omoimasu}
  \item \textquotedblleft{}I think it is dangerous\textquotedblright{} — \emph{abunai to omoimasu}
\end{itemize}

\begin{tcolorbox}[breakable,colback=teal!4,colframe=teal!35!black,title={Before you move on}]
\textquotedblleft{}I think he will come\textquotedblright{}? (\textbf{\ja{くると} \ja{おもいます}}.) \textquotedblleft{}I don't think he will come\textquotedblright{}? (\textbf{\ja{こないと} \ja{おもいます}}.) \textquotedblleft{}What do you think?\textquotedblright{} (\textbf{\ja{どう} \ja{おもいますか}}.) Which form goes before \textbf{\ja{と}}? (The plain one.)
\end{tcolorbox}
\section[\texorpdfstring{\ja{から} / \ja{ので} (kara / node)}{から / ので (kara / node)}]{\ja{から} / \ja{ので} (kara / node) — because}

\label{lesson:JA-C254-kara-node}



\begin{quote}
Say \textquotedblleft{}I don't think my friend will come\textquotedblright{}. (\emph{tomodachi wa konai to omoimasu}.) Then say \textquotedblleft{}cold\textquotedblright{} and \textquotedblleft{}a reason\textquotedblright{}. (\emph{samui}, \emph{riyū}.)
\end{quote}

\subsection*{You'll want to know: \ja{から} after a clause}
\begin{quote}
\textbf{\ja{さむいから、いきません。}}
\end{quote}

\begin{quote}
\emph{samui kara, ikimasen.} — \textquotedblleft{}It is cold, so I won't go.\textquotedblright{}
\end{quote}

\textbf{\ja{から}} (\emph{kara}) after a clause gives the reason, and the result follows it. The clause before it takes a plain form:

\begin{quote}
\textbf{\ja{じかんが} \ja{ないから、いそぎます。}}
\end{quote}

\begin{quote}
\emph{jikan ga nai kara, isogimasu.} — \textquotedblleft{}I have no time, so I will hurry.\textquotedblright{}
\end{quote}

\begin{quote}
\textbf{\ja{あめが} \ja{ふったから、いきませんでした。}}
\end{quote}

\begin{quote}
\emph{ame ga futta kara, ikimasen deshita.} — \textquotedblleft{}It rained, so I did not go.\textquotedblright{}
\end{quote}

\subsection*{You'll want to know: \ja{ので}, and \ja{りゆうは} …\ja{からです}}
\textbf{\ja{ので}} (\emph{node}) gives a reason too, more softly; it is the one to use when asking for something or excusing yourself:

\begin{quote}
\textbf{\ja{あたまが} \ja{いたいので、やすみます。}}
\end{quote}

\begin{quote}
\emph{atama ga itai node, yasumimasu.} — \textquotedblleft{}My head hurts, so I will rest.\textquotedblright{}
\end{quote}

When someone asks why, \textbf{\ja{りゆうは}} (\emph{riyū wa}), \textquotedblleft{}the reason is\textquotedblright{}, opens the answer, and \textbf{\ja{からです}} closes it:

\begin{quote}
\textbf{\ja{りゆうは} \ja{じかんが} \ja{ないからです。}}
\end{quote}

\begin{quote}
\emph{riyū wa jikan ga nai kara desu.} — \textquotedblleft{}The reason is that I have no time.\textquotedblright{}
\end{quote}

\subsection*{Your turn}
Say these aloud:

\begin{itemize}
\raggedright
  \item \textquotedblleft{}My head hurts, so I will rest\textquotedblright{} — \emph{atama ga itai node, yasumimasu}
  \item \textquotedblleft{}My stomach hurts, so I will rest\textquotedblright{} — \emph{onaka ga itai node, yasumimasu}
  \item \textquotedblleft{}My foot hurts, so I will rest\textquotedblright{} — \emph{ashi ga itai node, yasumimasu}
\end{itemize}

\begin{tcolorbox}[breakable,colback=teal!4,colframe=teal!35!black,title={Before you move on}]
Which comes first, the reason or the result? (The reason, with \textbf{\ja{から}} or \textbf{\ja{ので}}.) \textquotedblleft{}The reason is that I have no time\textquotedblright{}? (\textbf{\ja{りゆうは} \ja{じかんが} \ja{ないからです}}.)
\end{tcolorbox}
\section[\texorpdfstring{\ja{言}\ja{い}\ja{出}\ja{す} (iidasu)}{言い出す (iidasu)}]{\ja{言}\ja{い}\ja{出}\ja{す} (iidasu) — to suggest; to be the first to say}

\label{lesson:JA-C254-iidasu}



\begin{quote}
Before the new word: say \textquotedblleft{}to say\textquotedblright{} and \textquotedblleft{}to take out\textquotedblright{}. (\emph{iu}, \emph{dasu}.)
\end{quote}

\subsection*{You'll want to know: \ja{言}\ja{い}\ja{出}\ja{す}}
\textbf{\ja{言}\ja{い}\ja{出}\ja{す}} — \emph{iidasu} — \textquotedblleft{}to suggest\textquotedblright{}, \textquotedblleft{}to be the first to say\textquotedblright{}. Polite \textbf{\ja{言}\ja{い}\ja{出}\ja{します}}, \emph{iidashimasu}; a godan verb. Past (polite): \textbf{\ja{ちちが} \ja{りょこうを} \ja{言}\ja{い}\ja{出}\ja{しました}}, \emph{chichi ga ryokō o iidashimashita}.

\begin{quote}
\textbf{\ja{ちちが} \ja{りょこうを} \ja{言}\ja{い}\ja{出}\ja{しました。}}
\end{quote}

\begin{quote}
\emph{chichi ga ryokō o iidashimashita.} — \textquotedblleft{}My father suggested a trip.\textquotedblright{}
\end{quote}

It is \textbf{\ja{言}\ja{う}} with \textbf{\ja{だす}}, \textquotedblleft{}to put out\textquotedblright{}, after it: to put an idea out first. The first kanji is the \textbf{\ja{言}} you write. The second is new, read \textbf{\ja{だ}} (\emph{da}) here; the next lesson writes it.

\begin{cousinweb}
\textbf{\ja{言}\ja{い}\ja{出}\ja{す}} (\emph{iidasu}) is a native Japanese word (wago), inherited rather than borrowed from Chinese: the stem of \textbf{\ja{言}\ja{う}} (\emph{iu}), \textquotedblleft{}to say\textquotedblright{}, and \textbf{\ja{だす}} (\emph{dasu}), \textquotedblleft{}to put out\textquotedblright{}, each written with its kanji and read the native way. JMdict gives \textquotedblleft{}to begin to say\textquotedblright{} and \textquotedblleft{}to suggest (doing)\textquotedblright{} among its senses, and KANJIDIC2 gives \ja{い}.\ja{う} and \ja{だ}.\ja{す} as the two kanji's native readings (JMdict, KANJIDIC2).
\end{cousinweb}

\subsection*{Your turn}
Say these aloud:

\begin{itemize}
\raggedright
  \item \emph{iidasu}
  \item \textquotedblleft{}My father suggested a trip\textquotedblright{} — \emph{chichi ga ryokō o iidashimashita}
  \item \textquotedblleft{}My mother suggested a trip\textquotedblright{} — \emph{haha ga ryokō o iidashimashita}
\end{itemize}

\begin{tcolorbox}[breakable,colback=teal!4,colframe=teal!35!black,title={Before you move on}]
\textquotedblleft{}To suggest\textquotedblright{}? (\textbf{\ja{言}\ja{い}\ja{出}\ja{す}}, \emph{iidasu}.) Its polite past? (\textbf{\ja{言}\ja{い}\ja{出}\ja{しました}}.) How is the new kanji read here? (\textbf{\ja{だ}}, \emph{da}.)
\end{tcolorbox}
\section[\texorpdfstring{\ja{出} (da)}{出 (da)}]{\ja{出} — out, five strokes}

\label{lesson:JA-W254-da-kanji}



\begin{quote}
Say \emph{iidasu}, to suggest. Its first kanji is \textbf{\ja{言}}, which you write. Then call up \textbf{\ja{年}}: six strokes, the long stem last.
\end{quote}

\subsection*{Script — the shape}
\begin{quote}
\textbf{\ja{出}} — \emph{da}
\end{quote}

Its parts:

\begin{itemize}
\raggedright
  \item two open boxes, one stacked on the other
  \item one long centre stem through both
\end{itemize}

Five strokes, four lifts:

\begin{enumerate}
\raggedright
  \item draw the long centre stem from the top down
  \item lift, draw the upper left side down, then without lifting turn right along the base of the upper part
  \item lift, draw the upper right side down
  \item lift, draw the lower left side down, then without lifting turn right along the bottom
  \item lift, draw the lower right side down
\end{enumerate}

Unlike \textbf{\ja{年}}, the long stem comes first here, and the four sides hang from it. Here it is read \textbf{\ja{だ}} (\emph{da}), as in \textbf{\ja{言}\ja{い}\ja{出}\ja{す}}; this lesson asks only for that reading.

\subsection*{Your turn}
\begin{enumerate}
\raggedright
  \item Find \textbf{\ja{出}} in \textbf{\ja{言}\ja{い}\ja{出}\ja{す}} and say \emph{iidasu}.
  \item Copy \textbf{\ja{出}} three times, saying \emph{da} at each one.
  \item Find it once more, in \textbf{\ja{言}\ja{い}\ja{出}\ja{しました}} (\emph{iidashimashita}), and say the word.
\end{enumerate}

\begin{tcolorbox}[breakable,colback=teal!4,colframe=teal!35!black,title={Before you move on}]
How many strokes in all? (\textbf{Five}.)

Source: \href{https://github.com/KanjiVG/kanjivg/blob/master/kanji/051fa.svg}{KanjiVG stroke-order data}: KanjiVG, kanji/051fa.svg, kanji \ja{出} (U+51FA), 5 stroke paths kvg:051fa-s1 to s5 (Ulrich Apel and contributors, CC BY-SA 3.0). Only the order and direction are taken from it.
\end{tcolorbox}
\section[\texorpdfstring{\ja{いけん} (iken)}{いけん (iken)}]{\ja{いけん} (iken) — an opinion}

\label{lesson:JA-C254-iken}



\begin{quote}
Before the new word: say \textquotedblleft{}What do you think?\textquotedblright{} (\emph{dō omoimasu ka}.) Then \textquotedblleft{}my father suggested a trip\textquotedblright{}. (\emph{chichi ga ryokō o iidashimashita}.)
\end{quote}

\subsection*{You'll want to know: \ja{いけん}}
\textbf{\ja{いけん}} — \emph{iken} — \textquotedblleft{}an opinion\textquotedblright{}, what someone thinks.

\begin{quote}
\textbf{\ja{いけんは} \ja{ありますか。}}
\end{quote}

\begin{quote}
\emph{iken wa arimasu ka.} — \textquotedblleft{}Do you have an opinion?\textquotedblright{}
\end{quote}

\begin{quote}
\textbf{\ja{ちちの} \ja{いけんを} \ja{ききました。}}
\end{quote}

\begin{quote}
\emph{chichi no iken o kikimashita.} — \textquotedblleft{}I asked my father's opinion.\textquotedblright{}
\end{quote}

\begin{cousinweb}
\textbf{\ja{いけん}} (\emph{iken}) is Sino-Japanese (kango): two Chinese characters read with their on'yomi, readings borrowed from Chinese. They are \emph{i}, \textquotedblleft{}idea\textquotedblright{} (Middle Chinese \emph{'iH}) and \emph{ken}, \textquotedblleft{}see\textquotedblright{} (\emph{kenH}), which KANJIDIC2 also glosses \textquotedblleft{}opinion\textquotedblright{} (KANJIDIC2; Middle Chinese in Baxter's transcription). JMdict glosses the word \textquotedblleft{}opinion; view; comment\textquotedblright{} (JMdict).
\end{cousinweb}

\subsection*{Your turn}
Say these aloud:

\begin{itemize}
\raggedright
  \item \emph{iken}
  \item \textquotedblleft{}I asked my father's opinion\textquotedblright{} — \emph{chichi no iken o kikimashita}
  \item \textquotedblleft{}I asked my teacher's opinion\textquotedblright{} — \emph{sensei no iken o kikimashita}
\end{itemize}

\begin{tcolorbox}[breakable,colback=teal!4,colframe=teal!35!black,title={Before you move on}]
\textquotedblleft{}An opinion\textquotedblright{}? (\textbf{\ja{いけん}}, \emph{iken}.) \textquotedblleft{}Do you have an opinion?\textquotedblright{} (\textbf{\ja{いけんは} \ja{ありますか}}.)
\end{tcolorbox}
\section[(review)]{Review — what I think, and why}

\label{lesson:JA-R254-what-i-think-and-why}



\begin{quote}
Say \textquotedblleft{}I think it is good\textquotedblright{}. (\emph{ii to omoimasu}.) Then \textquotedblleft{}it is cold, so I won't go\textquotedblright{}. (\emph{samui kara, ikimasen}.)
\end{quote}

\subsection*{Your turn}
Say these aloud:

\begin{itemize}
\raggedright
  \item I think it will snow tomorrow.
  \item I don't think my friend will come.
  \item My head hurts, so I will rest.
  \item My father suggested a trip.
  \item Do you have an opinion?
\end{itemize}

\begin{tcolorbox}[breakable,colback=teal!4,colframe=teal!35!black,title={Before you move on}]
Which form goes before \textbf{\ja{と} \ja{おもいます}}? (The plain one.) Before \textbf{\ja{から}}? (The plain one too.) Which \textquotedblleft{}because\textquotedblright{} is softer? (\textbf{\ja{ので}}.) \textquotedblleft{}An opinion\textquotedblright{}? (\textbf{\ja{いけん}}.) \textquotedblleft{}To suggest\textquotedblright{}? (\textbf{\ja{言}\ja{い}\ja{出}\ja{す}}.)
\end{tcolorbox}
